\begin{ExerciseBlock}
\ExerciseHeading{01}{Schwerpunkt eines Bauteils}{rigid_bodies_composite_cylinder_centre_of_mass}
Ein Bauteil sei aus zwei kreisförmigen Vollzylindern
gleichen Materials zusammengesetzt, wie in der Skizze
dargestellt. Für die Durchmesser und Höhen gelte:
\[ d_2=d_1/2 \ \ \ \ \text{und} \ \ \ \ h_2 =2 h_1 .\]

Leiten Sie je eine Formel für die Schwerpunktkoordinaten
$x_S$ und $z_S$ her, wobei am Schluss in den Formeln nur
Daten von Zylinder 1, also $m_1$, $d_1$ und $h_1$ vorkommen.

\begin{center}
	\CachedTikzFigure{content/resources/chapters/rigid_bodies/tikz/composite_cylinder_centre_of_mass.tex}
\end{center}
\end{ExerciseBlock}
\UseExerciseSolution{rigid_bodies_composite_cylinder_centre_of_mass}
